\begin{abcliste}
\abc
$\theta_\mathrm{B} < \theta_\mathrm{C} < \theta_\mathrm{D} < \theta_\mathrm{A}$. Am stärksten geneigt ist A. Ob ein Körper umfällt, hängt aber nicht von der Neigung allein ab, sondern davon, wo sein Schwerpunkt bezüglich der Kante liegt (Teil b): \result{\text{A fällt nicht um.}}

\abc
Jeder Körper dreht sich um die Kante, auf der er steht. Die Gewichtskraft greift in S an und zeigt senkrecht nach unten. Verläuft die Lotrechte durch S innerhalb der Kante (über der Grundfläche), dreht die Gewichtskraft den Körper zurück; verläuft sie ausserhalb, fällt er um. Die Grenze ist erreicht, wenn S genau senkrecht über der Kante liegt, beim Kippwinkel
\[
 \tan\theta_\mathrm{k} = \frac{b/2}{h_\mathrm{S}}
\]
\begin{align}
 \theta_\mathrm{k,A} &= \kAF = \kA \approx \kAP{0}{3}\\
 \theta_\mathrm{k,B} &= \kBF = \kB \approx \kBP{0}{3}\\
 \theta_\mathrm{k,C} &= \kCF = \kC \approx \kCP{0}{3}\\
 \theta_\mathrm{k,D} &= \kDF = \kD \approx \kDP{0}{3}
\end{align}
Ein Körper fällt um, wenn $\theta > \theta_\mathrm{k}$:
\begin{itemize}
 \item A: $\tAO < \kAP{0}{3}$: fällt zurück.
 \item B: $\tBO > \kBP{0}{3}$: fällt um.
 \item C: $\tCO < \kCP{0}{3}$: fällt zurück.
 \item D: $\tDO > \kDP{0}{3}$: fällt um.
\end{itemize}
\[
 \result{\text{B und D fallen um, A und C fallen zurück.}}
\]
Zur Kontrolle der waagrechte Abstand von S zur Kante (positiv: S liegt noch über der Grundfläche, auf der Innenseite der Kante):
\begin{align}
 x &= \frac{b}{2}\cos\theta-h_\mathrm{S}\sin\theta\\
 x_\mathrm{A} &= \xA \approx \xAP{-2}{2}, \qquad x_\mathrm{B} = \xB \approx \xBP{-2}{2}\\
 x_\mathrm{C} &= \xC \approx \xCP{-2}{2}, \qquad x_\mathrm{D} = \xD \approx \xDP{-2}{2}
\end{align}
Der breite, unten beschwerte Körper A lässt sich am weitesten neigen; der schmale, hohe B fällt schon bei der kleinsten Neigung.
\end{abcliste}
